\conclusion

В ходе выпускной квалификационной работы была разработана система обнаружения аномалий в Kubernetes-кластере на основе методов нейронных сетей, предназначенная для анализа audit-логов и выявления отклонений от нормального поведения в контейнерной среде.

В исследовательской части были рассмотрены основные подходы к обнаружению аномалий в информационных системах, а также проанализированы особенности обеспечения безопасности в Kubernetes-кластерах. Был выполнен обзор статических методов, алгоритмов машинного обучения и нейросетевых подходов к анализу последовательностей событий. Установлено, что наибольший потенциал для решения поставленной задачи имеют методы глубокого обучения, способные учитывать временную структуру данных и выявлять сложные зависимости в журналируемых событиях. Отдельное внимание было уделено механизму audit-логирования Kubernetes, поскольку audit-логи API-сервера являются наиболее информативным источником данных для обнаружения аномалий в действиях пользователей, сервисных аккаунтов и компонентов кластера.

В конструкторской части была обоснована постановка задачи обнаружения аномалий на основе анализа последовательностей audit-событий. В условиях отсутствия достаточного объёма размеченных данных для всех возможных аномальных сценариев наиболее целесообразным был признан подход на основе автоэнкодеров, позволяющий обучать модель на данных, описывающих нормальное поведение системы. В рамках работы были рассмотрены и сравнены архитектуры AE, RNN-AE и LSTM-AE, а также определены входные данные, критерии обучения и метрики качества. Проведённый анализ показал, что использование рекуррентных нейросетевых моделей, в частности LSTM-автоэнкодера, является наиболее эффективным для выявления аномалий в последовательностях audit-логов. По результатам экспериментов именно модель LSTM-AE продемонстрировала наилучшие значения качества, обеспечив $AUC_{PR} = 0{,}99962$, precision = 0{,}93097, recall = 1{,}00000 и F1-мера = 0{,}96425.

В технологической части было разработано программное средство, реализующее полный цикл обработки audit-логов: от получения и предобработки данных до обучения модели и применения её для выявления аномалий. При разработке использовались современные программные инструменты и библиотечные средства, обеспечивающие надёжность, производительность и удобство сопровождения системы. Применение Python в качестве основного языка программирования, фреймворка FastAPI для серверной части, Uvicorn в качестве ASGI-сервера, PostgreSQL и SQLAlchemy для работы с данными, Celery для организации распределённых вычислений, а также PyTorch и scikit-learn для реализации и оценки нейросетевой модели позволило создать масштабируемое и технологически целостное решение. Разработанное приложение обеспечивает автоматизированный анализ журналов аудита и своевременное обнаружение аномальных событий в Kubernetes-кластере.

В организационно-экономической части была выполнена оценка трудоёмкости, затрат на разработку и реализацию программного продукта, а также определена его экономическая целесообразность. Результаты расчётов показали, что разработка проекта требует значительных, но обоснованных ресурсов, а построенный календарный план и оценка затрат подтверждают реалистичность выполнения работ в установленные сроки. Проведённый анализ свидетельствует о том, что разработка данного решения является экономически оправданной и востребованной на рынке средств обеспечения информационной безопасности.

В правовой части был проведён анализ нормативных правовых актов Российской Федерации, регулирующих вопросы информации, персональных данных и защиты информации. Установлено, что разрабатываемое программное средство соответствует требованиям действующего законодательства и может применяться в информационных системах при соблюдении необходимых организационных и технических мер защиты. Это подтверждает правомерность и практическую применимость предложенного решения.

Таким образом, в ходе выполнения выпускной квалификационной работы была достигнута поставленная цель, а именно разработано средство обнаружения аномалий в Kubernetes-кластере на основе методов нейронных сетей, обеспечивающее высокую точность выявления аномальных событий. Полученные результаты подтверждают эффективность выбранного подхода и создают основу для дальнейшего развития системы, расширения функциональности и её интеграции в практические сценарии мониторинга и обеспечения безопасности контейнерных инфраструктур.